\BourbakiExerciseGroup

\begin{enumerate}
\def\labelenumi{\arabic{enumi})}
\item
  Let A be a not necessarily commutative ring, not zero and without divisors of zero (I, p.~98). Show that if there exists an integer \(m \geqslant 1\) such that \(m\mathbf{A} = 0\) , then the set of all integers having this property is an ideal \(p\mathbf{Z}\) , where \(p\) is a prime number, and A is a ring of characteristic \(p\) .
\item
  Let \(m, n\) be integers such that \(0 \leqslant n \leqslant m\) . Let \(p\) be a prime number and let \(m = \alpha_0 + \alpha_1p + \cdots + \alpha_rp^r\) and \(n = \beta_0 + \beta_1p + \cdots + \beta_rp^r\) be the expansions of \(m\) and \(n\) in base \(p\) , where we thus have \(0 \leqslant \alpha_i < p\) , \(0 \leqslant \beta_i < p\) for \(0 \leqslant i \leqslant r\) ; suppose that \(\alpha_r \neq 0\) . Show that \(\binom{m}{n} \equiv 0 (\bmod p)\) if \(\beta_i > \alpha_i\) for at least one index \(i\) . If on the other hand, \(\beta_i \leqslant \alpha_i\) for \(0 \leqslant i \leqslant r\) , then
\end{enumerate}

\[
\binom {m} {n} \equiv \binom {\alpha_ {0}} {\beta_ {0}} \dots \binom {\alpha_ {r}} {\beta_ {r}} \pmod {p}  .
\]

(In \((\mathbf{Z} / p\mathbf{Z})[\mathbf{X}]\) we have \((1 + \mathbf{X})^m = (1 + \mathbf{X})^{\alpha_0}(1 + \mathbf{X}^p)^{\alpha_1}\dots (1 + \mathbf{X}^{p^r})^{\alpha_r})\)

\begin{enumerate}
\def\labelenumi{\arabic{enumi})}
\setcounter{enumi}{2}
\tightlist
\item
  Let \(\mathbf{A}\) be a commutative ring. For every polynomial \(f \in \mathbf{A}[\mathbf{X}]\) we write in the polynomial ring \(\mathbf{A}[\mathbf{X}, \mathbf{Y}]\) of polynomials in two indeterminates over \(\mathbf{A}\) ,
\end{enumerate}

\[
f (\mathbf {X} + \mathbf {Y}) = \sum_ {m = 0} ^ {\infty} \Delta_ {m} f (\mathbf {X}) \mathbf {Y} ^ {m},
\]

where \(\Delta_{m}\) is thus an A-linear mapping of A{[}X{]} into itself.\\
a) If Z is an indeterminate and \(\tau_{Z}\) the A-linear mapping of A{[}X{]} into A{[}X, Z{]} such that \(\tau_{Z}f = f(X + Z)\) , show that \(\tau_{Z} \circ \Delta_{m} = \Delta_{m} \circ \tau_{Z}\) (where on the right \(\Delta_{m}\) is the mapping defined above in the ring B{[}X{]}, where B = A{[}Z{]}). Deduce that in A{[}X{]} we have

\[
\Delta_ {m} \Delta_ {n} = \Delta_ {n} \Delta_ {m} = \binom {m + n} {m} \Delta_ {m + n}
\]

for integers \(m, n \geqslant 0\) (cf.~IV, p.~8).

\begin{enumerate}
\def\labelenumi{\alph{enumi})}
\setcounter{enumi}{1}
\tightlist
\item
  Suppose that A is a ring of characteristic p \textgreater{} 0 and for each integer \(k \geqslant 0\) put \(D_{k} = \Delta_{p^{k}}\) . Show that if \(f \in A[X]\) and \(g \in A[X^{p^{k}}]\) then \(D_{k}(fg) = g \cdot D_{k}f + f \cdot D_{k}g\) .
\end{enumerate}

Every polynomial \(f \in \mathbf{A}[\mathbf{X}]\) can be uniquely written in the form \(f(\mathbf{X}) = \sum_{m=0}^{\infty} \mathbf{X}^{mp^k} g_m(\mathbf{X})\) where \(\deg(g_m) < p^k\) ; show that

\[
\mathrm{D} _ {k} f (\mathbf {X}) = \sum_ {m = 0} ^ {\infty} m \mathbf {X} ^ {(m - 1) p ^ {k}} g _ {m} (\mathbf {X})
\]

and deduce that \(D_{k}^{p}=0\) .

\begin{enumerate}
\def\labelenumi{\alph{enumi})}
\setcounter{enumi}{2}
\tightlist
\item
  For every integer m \textgreater{} 0 let \(m = \alpha_{0} + \alpha_{1}p + \cdots + \alpha_{k}p^{k}\) be the expansion of m in base p (Set Theory, III, p.~177), where thus \(0 \leqslant \alpha_{j} < p\) for \(0 \leqslant j \leqslant k\) ; suppose that \(\alpha_{k} \neq 0\) . Show that
\end{enumerate}

\[
\Delta_ {m} = \left(\frac {1}{\alpha_ {0} !} D _ {0} ^ {\alpha_ {0}}\right) \left(\frac {1}{\alpha_ {1} !} D _ {1} ^ {\alpha_ {1}}\right) \dots \left(\frac {1}{\alpha_ {k} !} D _ {k} ^ {\alpha_ {k}}\right).
\]

(If \(\Delta_m'\) is the right-hand side of this relation, note that if we introduce the expansion \(f(\mathbf{X} + \mathbf{Y}) = \sum_{q=0}^{\infty} f_q(\mathbf{Y}) \mathbf{X}^q\) , we have \(\Delta_m' f(\mathbf{X} + \mathbf{Y}) = \sum_{q=0}^{\infty} f_q(\mathbf{Y}) \Delta_m'(\mathbf{X}^q)\) and deduce that this may be written \(\Delta_m' f(\mathbf{X} + \mathbf{Y}) = f_m(\mathbf{X}) + \mathbf{Y} g(\mathbf{X}, \mathbf{Y})\) , where \(g\) is a polynomial.)

\begin{enumerate}
\def\labelenumi{\arabic{enumi})}
\setcounter{enumi}{3}
\tightlist
\item
  Let \(\mathbf{G}\) be a commutative group, written additively, and let \(p\) be a prime number. Denote by \(\mathbf{G}\left[\frac{1}{p}\right]\) the commutative group \(\mathbf{G} \otimes_{\mathbf{Z}} \mathbf{Z}\left[\frac{1}{p}\right]\) . Show that if \(\mathbf{K}\) is a field of characteristic \(p\) , then the perfect closure of the group algebra \(\mathbf{K}[\mathbf{G}]\) is the group algebra \(\mathbf{K}^{p^{-\infty}}\left[\mathbf{G}\left[\frac{1}{p}\right]\right]\) .
\end{enumerate}

* 5) Let \(\mathfrak{P}\) be the set of prime numbers and put \(\mathbf{A} = \prod_{p\in \mathfrak{P}}\mathbf{F}_p\) . For every filter \(\mathfrak{F}\) on \(\mathfrak{P}\) (Gen.

Top., I, p.~57) denote by \(m_{\mathfrak{F}} \subset \mathbf{A}\) the set of \(u = (u_p)_{p \in \mathfrak{P}}\) such that the set of \(p \in \mathfrak{P}\) with \(u_p = 0\) is a member of \(\mathfrak{F}\) . Show that \(m_{\mathfrak{F}}\) is an ideal and that we thus obtain a bijection between the ideals of \(\mathbf{A}\) distinct from \(\mathbf{A}\) and the filters on \(\mathfrak{P}\) . Let \(\mathfrak{F}\) be an ultrafilter (Gen.~Top., I, p.~60) on \(\mathfrak{P}\) which does not possess a limit in the discrete topology. Show that \(\mathbf{A}/m_{\mathfrak{F}}\) is a field of characteristic 0, isomorphic to a subfield of \(\mathbf{C}\) . *

